\documentclass[11pt,a4paper]{article}
% Standalone: no \input, no BibTeX, no external files.
\usepackage[utf8]{inputenc}
\usepackage[T1]{fontenc}
\usepackage{lmodern}
\usepackage{amsmath,amssymb,amsthm}
\setlength{\emergencystretch}{4em}
\setlength{\hfuzz}{2pt}
\usepackage{geometry}
\geometry{margin=2.5cm}
\usepackage[numbers,sort&compress]{natbib}
\usepackage[hidelinks]{hyperref}
\theoremstyle{plain}
\newtheorem{theorem}{Theorem}[section]
\newtheorem{lemma}[theorem]{Lemma}
\newtheorem{proposition}[theorem]{Proposition}
\newtheorem{corollary}[theorem]{Corollary}
\theoremstyle{definition}
\newtheorem{definition}[theorem]{Definition}
\newtheorem{example}[theorem]{Example}
\theoremstyle{remark}
\newtheorem{remark}[theorem]{Remark}
\DeclareMathOperator{\Tr}{Tr}
\DeclareMathOperator{\TV}{TV}
\DeclareMathOperator{\spec}{spec}
\newcommand{\norm}[1]{\left\lVert #1\right\rVert}
\newcommand{\ip}[2]{\left\langle #1,#2\right\rangle}
\newcommand{\C}{\mathbb{C}}
\newcommand{\R}{\mathbb{R}}

\title{Energy Certificates and Conditional Kinetic Bounds for
Finite-Basis Surface-Reaction Models}
\author{}
\date{}
\providecommand{\doi}[1]{\href{https://doi.org/#1}{\texttt{doi:#1}}}
\providecommand{\path}[1]{\texttt{#1}}

\begin{document}
\maketitle

\begin{abstract}
A computed electronic energy becomes an error certificate only when both
sides of an interval are justified for a specified Hamiltonian. This research
note gives a self-contained chain of conditional bounds: a finite-dimensional
block decomposition yields a ground-energy enclosure; interval subtraction
yields an electronic barrier enclosure; separately bounded thermodynamic and
geometric corrections allow propagation through an assumed transition-state
rate formula; and a finite-state Markov generator perturbation bound controls
state probabilities and bounded observables. The block bound accepts lower
bounds on both diagonal blocks and an upper bound on the coupling norm. Its
proof exposes the principal computational difficulty: certifying the omitted
space. We prove the direction of reduced-density-matrix relaxation bounds,
distinguish dual certificates from approximate primal values, and explain why
a small eigenvector residual need not identify a ground state. A two-level
pairing model and a three-vertex conflict graph provide exact algebraic checks
and a counterexample to a proposed topological occupancy rule. These are
elementary mathematical results assembled for certificate design; no novelty
is claimed for the inequalities. No molecular or surface benchmark, uniform
Born--Oppenheimer error estimate, universal polynomial-time algorithm, or
experimental growth prediction is established. The principal open task is to
construct useful, inexpensive, rigorously evaluated inputs for one fully
specified surface-reaction model.
\end{abstract}

\section{Purpose and scope}
\label{sec:intro}
The useful question for a reduced electronic calculation is whether it can
exclude an error large enough to change a chemical conclusion. A trial state
can give a low energy without supplying such an exclusion. Similarly, close
agreement between implementations can validate shared algebra while leaving
the physical model untested. We separate three objects: a mathematical target,
a certificate for that target, and the additional assumptions relating the
target to an experiment.

Our target is the lowest eigenvalue of a specified finite Hermitian matrix,
usually an electronic Hamiltonian in a fixed orbital basis, particle-number
sector, geometry, and boundary condition. Its energy interval does not by
itself bound the continuum electronic problem, a free energy, a transition
rate, or growth per cycle. Each passage to a different target requires an
explicit model or an additional error bound.

Worst-case complexity results for electronic structure in a fixed basis
\citep{OGorman2022} and for fermionic $N$-representability \citep{Liu2007}
prevent an inference of general efficiency from the notation of chemistry or
from replacing a wavefunction by reduced data. They do not say that every
chemical instance is difficult. Our results make no average-case complexity
claim. They identify sufficient information for a sound conclusion and make
visible where obtaining that information may be expensive.

Ab initio approaches to area-selective ALD motivate the need to separate
surface-model choices, electronic accuracy, and kinetic assumptions
\citep{Pieck2025}. The intended applications include elementary steps involving
$\mathrm{Al_2O_3}/\mathrm{MoCl_5}$ and $\mathrm{CaF_2}$ crystal deposition, as
specified by the user. These are research targets, not descriptions of
completed calculations. No surface structures, molecular integrals, independent
DFT labels, precursor pair, or measured kinetic data for either target are
provided here; in particular no reaction stoichiometry for the
$\mathrm{CaF_2}$ process has been fixed.

\section{Finite-dimensional setting and block energy certificates}
\label{sec:energy}
\subsection{The target and the meaning of a certificate}
Let $\mathcal H$ be a nonzero finite-dimensional complex Hilbert space and
let $H=H^*$ on $\mathcal H$. We use inner products conjugate-linear in their
first argument and the induced operator norm. Define
\[
 E_0(H)=\min_{\norm{v}=1}\ip{v}{Hv}=\lambda_{\min}(H).
\]
The equality follows by expanding $v$ in an orthonormal eigenbasis: its
expectation is a convex combination of eigenvalues, and a lowest eigenvector
attains the minimum. An energy certificate is a pair of justified real
numbers $L,U$ with $L\le E_0(H)\le U$. A floating-point estimate of an
endpoint is not automatically a justified endpoint.

For fermionic applications, take $K$ orthonormal spin orbitals and the
$N$-particle sector $\bigwedge^N\C^K$, where $0\le N\le K$.
Its dimension is $\binom K N$. One possible convention is
\begin{equation}
 H=E_{\mathrm{nuc}}I+\sum_{pq}h_{pq}a_p^\dagger a_q
   +\frac12\sum_{pqrs}v_{pqrs}a_p^\dagger a_q^\dagger a_s a_r,
 \label{eq:electronic}
\end{equation}
with coefficients chosen so that $H$ is Hermitian. The operators obey
$\{a_p,a_q^\dagger\}=\delta_{pq}I$ and
$\{a_p,a_q\}=\{a_p^\dagger,a_q^\dagger\}=0$.
The coefficient convention, energy units, nuclear term, and orbital basis
must be supplied with a numerical instance. The theorems below do not assume
any particular form of these coefficients.

If one additionally restricts total spin or another symmetry, the target is
the ground energy in that sector. A certificate for one sector is not a
global ground-energy certificate unless the other sectors have been excluded
or bounded. Across a complete direct-sum decomposition, the global ground
energy is the minimum of the sector ground energies.

\subsection{A bound using certified block data}
Choose an orthogonal decomposition
$\mathcal H=\mathcal H_A\oplus\mathcal H_C$ with both summands nonzero:
\begin{equation}
 H=\begin{pmatrix} A&B\\B^*&C\end{pmatrix}.
 \label{eq:block}
\end{equation}
Thus $B:\mathcal H_C\to\mathcal H_A$. The active space may be small,
but the complement comprises every other vector in the stated target space.

\begin{theorem}[Block enclosure from one-sided data]
\label{thm:block}
Suppose real numbers $\underline a,\underline c,\overline b$ satisfy
\[
 A\succeq\underline a I,\qquad C\succeq\underline c I,
 \qquad \norm B\le\overline b,\qquad \overline b\ge0.
\]
Let $x_0\in\mathcal H_A$ be normalized and let
$u=\ip{x_0}{Ax_0}$, or more generally let $u$ be a certified upper bound
on this expectation. Then
\begin{equation}
 \ell:=\frac{\underline a+\underline c
       -\sqrt{(\underline c-\underline a)^2+4\overline b^2}}2
 \ \le\ E_0(H)\ \le\ u.
 \label{eq:block-enclosure}
\end{equation}
No ordering of $\underline a$ and $\underline c$ is required.
\end{theorem}
\begin{proof}
Write a normalized vector as $v=(x,y)$, and put
$s=\norm x$, $t=\norm y$. Then $s^2+t^2=1$. The hypotheses and
Cauchy--Schwarz give
\begin{align*}
 \ip{v}{Hv}
 &=\ip{x}{Ax}+\ip{y}{Cy}+2\operatorname{Re}\ip{x}{By}\\
 &\ge\underline a s^2+\underline c t^2-2\overline b st.
\end{align*}
The final expression is the quadratic form of
$M=\left(\begin{smallmatrix}\underline a&-\overline b\\
-\overline b&\underline c\end{smallmatrix}\right)$
on the unit vector $(s,t)$. Its characteristic polynomial is
$(\underline a-z)(\underline c-z)-\overline b^2$; the smaller root is
$\ell$. Expanding $(s,t)$ in an orthonormal eigenbasis of this real
symmetric matrix gives $(s,t)M(s,t)^T\ge\ell$. Minimizing over $v$
proves the lower bound. The vector $(x_0,0)$ is a normalized trial vector
for $H$, with expectation $\ip{x_0}{Ax_0}\le u$. The variational
principle proves the upper bound.
\end{proof}

\begin{corollary}[Separated complement]
\label{cor:gap}
Suppose $a=\lambda_{\min}(A)$ is known exactly,
$C\succeq cI$, $\norm B\le b$, $b\ge0$, and $c>a$. Then
\begin{equation}
 0\le a-E_0(H)
 \le\frac{\sqrt{(c-a)^2+4b^2}-(c-a)}2
 =\frac{2b^2}{\sqrt{(c-a)^2+4b^2}+(c-a)}
 \le\frac{b^2}{c-a}.
 \label{eq:gap-bound}
\end{equation}
If $a$ is only replaced by a lower bound $\underline a$, the assertion
$0\le\underline a-E_0(H)$ need not hold; the applicable enclosure remains
\eqref{eq:block-enclosure} with an independently justified upper endpoint.
\end{corollary}
\begin{proof}
A normalized lowest eigenvector of $A$ gives $E_0(H)\le a$.
Theorem~\ref{thm:block} with input $(a,c,b)$ gives the middle bound by
subtracting its lower endpoint from $a$. Rationalizing the numerator yields
the identity. Since $c-a>0$, the denominator is at least $2(c-a)$,
which proves the last inequality, including $b=0$.
\end{proof}

\begin{proposition}[Sharpness of the block estimate]
\label{prop:sharp}
For any real $a,c$ and $b\ge0$, the scalar-block matrix
$H=\left(\begin{smallmatrix}a&-b\\-b&c\end{smallmatrix}\right)$
attains the lower endpoint of Theorem~\ref{thm:block} with exact inputs.
For fixed $c-a>0$, the coefficient one in the final bound
$b^2/(c-a)$ cannot be reduced uniformly for all $b>0$.
\end{proposition}
\begin{proof}
The characteristic polynomial in the preceding proof is now precisely that
of $H$, so its smaller root is $E_0(H)$. Writing $d=c-a>0$, the ratio
of the exact shift to $b^2/d$ is
$2d/(\sqrt{d^2+4b^2}+d)$, which tends to one as $b$ decreases to zero.
Any proposed smaller constant fails for sufficiently small positive $b$.
\end{proof}

When $B=0$, the lower endpoint is $\min(\underline a,\underline c)$.
A low state in the complement remains visible. The square-root bound also
remains valid when $c=a$ or $c<a$, where Corollary~\ref{cor:gap} cannot
be used.

\subsection{What the certificate costs}
The theorem does not compute its hypotheses. A ground-energy approximation
to $C$ obtained by a variational method is an upper bound on
$\lambda_{\min}(C)$ and cannot serve as $\underline c$.
Diagonalizing $C$ may cost essentially as much as the original calculation.
A bound on a one-particle orbital gap is not, without a further argument,
a lower bound on every many-electron energy in the omitted space.

Crude rigorous bounds can nevertheless be useful. For example, if an explicit
Hermitian matrix $C=(C_{ij})$ is available, then
\begin{equation}
 C\succeq
 \left[\min_i\left(C_{ii}-\sum_{j\ne i}|C_{ij}|\right)\right]I.
 \label{eq:row-bound}
\end{equation}
Indeed, each off-diagonal contribution obeys
$2\operatorname{Re}(\overline y_i C_{ij}y_j)
\ge-|C_{ij}|(|y_i|^2+|y_j|^2)$; summing over $i<j$ proves the
quadratic-form inequality. Likewise $\norm B\le\norm B_F$, since
Cauchy--Schwarz applied to each row gives
$\norm{By}^2\le(\sum_{ij}|B_{ij}|^2)\norm y^2$.
These estimates may be inexpensive for a sparse represented matrix but very
loose, and even reading all entries may be infeasible in a large Fock space.
A practical method must state how the matrix or operator is accessed.

All scalar evaluation must preserve the one-sided inequalities. Exact
arithmetic, outward-rounded intervals, or independently verified error bounds
can do this. An ordinary dense eigensolver and a rounded square root provide
numerical checks, not a rounding certificate. A cancellation-resistant identity
is
\[
 \ell=\min(\underline a,\underline c)
 -\frac{2\overline b^2}
 {\sqrt{(\underline c-\underline a)^2+4\overline b^2}
       +|\underline c-\underline a|},
\]
with the correction defined as zero when $\overline b=0$. Stability of this
formula still does not replace outward rounding.

\section{Variational states and reduced-density-matrix relaxations}
\label{sec:rdm}
\subsection{Upper bounds and the limits of residuals}
For a normalized trial state $\psi\in\mathcal H$, the expectation
$u=\ip{\psi}{H\psi}$ satisfies $E_0(H)\le u$. The same holds for
$\Tr(H\Gamma)$ when $\Gamma\succeq0$ and $\Tr\Gamma=1$: in an
eigenbasis of $H$, the diagonal entries of $\Gamma$ are nonnegative
weights summing to one. These statements require neither an optimized state
nor an exact representation of the ground state. They do require an
accurately bounded expectation for the stated $H$.

\begin{proposition}[What a residual proves]
\label{prop:residual}
For a normalized $\psi$, a real number $\mu$, and
$r=\norm{(H-\mu I)\psi}$, at least one eigenvalue of $H$ lies in
$[\mu-r,\mu+r]$. This assertion does not identify the ground eigenvalue.
\end{proposition}
\begin{proof}
Let $H\phi_j=\lambda_j\phi_j$ be an orthonormal eigenbasis, and write
$\psi=\sum_j c_j\phi_j$. Then
\[
 r^2=\sum_j|c_j|^2(\lambda_j-\mu)^2
 \ge\min_j|\lambda_j-\mu|^2\sum_j|c_j|^2
 =\min_j|\lambda_j-\mu|^2.
\]
Taking square roots proves the spectral-distance statement. For the claimed
limitation, choose $H=\operatorname{diag}(-2,1)$ and $\psi=(0,1)^T$.
Its Rayleigh quotient is $\mu=1$ and its residual is zero, although
$E_0(H)=-2$. Thus even a vanishing residual can belong to an excited state.
\end{proof}

An a posteriori theorem using a residual to bound the ground energy must
supply additional information that selects the ground branch, such as an
appropriate certified spectral separation. A small residual alone does not
provide the missing lower bound on the complement in
Theorem~\ref{thm:block}.

\subsection{Outer relaxations: the direction of the inequality}
For a density operator $\Gamma$ supported in the chosen $N$-particle
space, define
\[
 \gamma_{pq}=\Tr(\Gamma a_p^\dagger a_q),\qquad
 D_{pq,rs}=\Tr(\Gamma a_p^\dagger a_q^\dagger a_s a_r).
\]
Let $\mathcal D_N$ be the set of all pairs $(\gamma,D)$ so obtained.
Keeping both matrices avoids any special case in a contraction formula at
$N=0$ or $N=1$. The electronic energy is the affine functional
\[
 \mathcal E(\gamma,D)=E_{\mathrm{nuc}}+
 \sum_{pq}h_{pq}\gamma_{pq}+\frac12\sum_{pqrs}v_{pqrs}D_{pq,rs}.
\]
By the density-operator variational principle,
$E_0(H)=\min_{d\in\mathcal D_N}\mathcal E(d)$.
The use of necessary positivity constraints to relax this representability
condition is part of the variational RDM approach \citep{Mazziotti2012}.

\begin{proposition}[Relaxation and restriction]
\label{prop:relaxation}
Suppose
$\mathcal D_N\subseteq\mathcal R_{\mathrm{tight}}
\subseteq\mathcal R_{\mathrm{loose}}$ and define
$e_{\mathcal R}=\inf_{d\in\mathcal R}\mathcal E(d)$.
Then
\begin{equation}
 e_{\mathrm{loose}}\le e_{\mathrm{tight}}\le E_0(H)\le u
 \label{eq:relaxation}
\end{equation}
for every valid trial-state upper bound $u$.
In contrast, minimization over a nonempty subset of physical states gives an
upper bound on $E_0(H)$.
\end{proposition}
\begin{proof}
Every point in the smaller feasible set is also in the larger set, so the
infimum over the larger set cannot be greater. Applying this observation to
the two inclusions proves the first two inequalities. The trial-state
variational principle proves the third. For a restriction to physical
states, each feasible energy is at least $E_0(H)$, so its infimum is also
at least $E_0(H)$.
\end{proof}

Here an infimum may be $-\infty$, in which case it is a valid but useless
lower bound. Additional valid outer-relaxation constraints can only raise
the optimum. Percentages of a correlation energy, relative errors, and
absolute errors in energy units are distinct quantities. None can be
converted into another without the needed reference energy and convention.
In particular, no universal accuracy guarantee for a fixed RDM hierarchy
level follows from \eqref{eq:relaxation}.

\subsection{A dual objective is a certificate only after feasibility is checked}
Suppose an outer relaxation has the semidefinite representation
\begin{equation}
 e_{\mathrm{relax}}=
 \inf\{\langle F,X\rangle:\mathcal A(X)=d,\ X\succeq0\},
 \label{eq:sdp-primal}
\end{equation}
where $\langle F,X\rangle=\operatorname{Re}\Tr(F^*X)$ and
$\mathcal A$ maps Hermitian matrices to a real vector space. An additive
energy constant can be kept on both primal and dual sides. Define the
adjoint by $\langle\mathcal A^*y,X\rangle=y\cdot\mathcal A(X)$.

\begin{proposition}[Weak duality and a certified residual correction]
\label{prop:dual}
If $F-\mathcal A^*y\succeq0$, then
\[
 d\cdot y\le e_{\mathrm{relax}}\le E_0(H).
\]
More generally, suppose all primal-feasible $X$ have
$\Tr X\le\tau$ and a verified $\delta\ge0$ satisfies
$F-\mathcal A^*y\succeq-\delta I$. Then
$d\cdot y-\delta\tau\le e_{\mathrm{relax}}$.
\end{proposition}
\begin{proof}
For any feasible $X$, put $S=F-\mathcal A^*y$. Then
\[
 \langle F,X\rangle-d\cdot y=\Tr(SX).
\]
If $S,X\succeq0$, this trace is nonnegative: it equals
$\Tr(X^{1/2}SX^{1/2})$, the trace of a positive semidefinite matrix.
Taking the infimum proves weak duality. In the second case,
$S+\delta I\succeq0$, so
$\Tr(SX)\ge-\delta\Tr X\ge-\delta\tau$; taking the same infimum
proves the corrected bound. The comparison with $E_0(H)$ follows from
the assumed outer inclusion.
\end{proof}

No strong-duality assumption is needed for this one-sided certificate.
A solver's approximate feasible primal point supplies an upper bound on the
relaxation optimum, not a lower bound on $E_0(H)$; if it is not
$N$-representable it supplies no trial-state upper bound either. Even in
$\min_{0\le x\le1}x=0$, the feasible objective $x=1/2$ is not a lower
bound on the minimum. A raw approximate dual objective is also insufficient
when its slack is not known to be positive semidefinite. Exact feasibility
or a rigorously bounded correction, such as Proposition~\ref{prop:dual},
is essential. Numerical solver residuals with unverified tolerances do not
meet this requirement by themselves.

One can combine independent lower bounds by taking their maximum. If
$L_{\mathrm{dual}}$ and the block endpoint $\ell$ are both justified for
the same Hamiltonian and sector, then
$\max(L_{\mathrm{dual}},\ell)\le E_0(H)\le u$. Applying a relaxation
to a complementary block is permissible only when its feasible set is proved
to contain every state in that block.

\subsection{Why algebraic constraints need physical inclusion proofs}
A proposed new constraint must hold for every physical $N$-particle state.
Testing it on selected states cannot establish this inclusion. Two exact
examples show why apparently stronger formulas may fail.

\begin{example}[The 2-RDM does not determine the 3-RDM]
\label{ex:three-rdm}
In six ordered spin orbitals, use the determinant convention
$|ijk\rangle=a_i^\dagger a_j^\dagger a_k^\dagger|0\rangle$ for
$i<j<k$. The normalized states
\[
 |\psi_\pm\rangle=(|123\rangle\pm|456\rangle)/\sqrt2
\]
have identical 1- and 2-RDMs. Indeed, a monomial with at most two annihilators
can remove at most two of the three occupied orbitals of $|456\rangle$;
with at most two creators it cannot turn that determinant into $|123\rangle$.
The reverse off-diagonal matrix element also vanishes. Thus every one- and
two-body expectation contains the same diagonal terms in the two states and
no sign-sensitive cross term. On the other hand,
\[
 \langle\psi_\pm|
 a_1^\dagger a_2^\dagger a_3^\dagger a_6a_5a_4
 |\psi_\pm\rangle=\pm\tfrac12,
\]
because the operator maps $|456\rangle$ to $|123\rangle$ with positive
sign and annihilates $|123\rangle$. Therefore no state-independent function
of the 2-RDM can reconstruct the entire 3-RDM.
\end{example}

\begin{example}[An invalid proposed necessary inequality]
For $|\psi\rangle=(|12\rangle+|34\rangle)/\sqrt2$, one has
$\langle n_1\rangle=\langle n_2\rangle=1/2$ and
$D_{12,12}=\langle n_1n_2\rangle=1/2$, by evaluating each determinant.
Consequently the proposed inequality
$D_{12,12}^2\le\langle n_1\rangle\langle n_2\rangle D_{12,12}$
would require $1/4\le1/8$. It excludes a physical state and cannot define a
sound outer relaxation of all $N$-representable data.
\end{example}

A valid operator sum-of-squares certificate is different from testing a
polynomial at random points. If, on the specified particle sector, one proves
$H-LI=\sum_\alpha T_\alpha^*T_\alpha$ (possibly modulo identities
that vanish on that sector), then every normalized expectation is at least
$L$, since each term contributes $\norm{T_\alpha\psi}^2\ge0$.
A practical SOS relaxation must specify its operator basis, the CAR and
particle-number identities, and a verified Gram or dual certificate. No such
new scalable SOS hierarchy or universal convergence rate is established here.

\section{Trial-state compression and topological descriptors}
\label{sec:tensor}
\subsection{Compression supplies trial states, not omitted-space bounds}
An open-boundary matrix product state on $L$ sites with local dimension $q$
has coefficients
\[
 \psi_{i_1\cdots i_L}
 =A^{[1]}_{i_1}A^{[2]}_{i_2}\cdots A^{[L]}_{i_L},
 \qquad A^{[j]}_{i_j}\in\C^{D_{j-1}\times D_j},\quad D_0=D_L=1.
\]
The displayed representation contains
$\sum_{j=1}^L qD_{j-1}D_j$ scalar entries. If $q$ is fixed and every
$D_j\le D$, this number is at most $LqD^2$. This is a storage bound
for a supplied representation. It does not bound the bond dimension required
for a given energy accuracy, the time to find suitable tensors, or the cost
of contracting a general Hamiltonian representation.

A normalized state in the intended particle and symmetry sector gives a
variational upper bound once its expectation is rigorously evaluated. Orbital
ordering and local tensor optimization can be used to search for that state;
the certificate is the resulting energy enclosure, not the optimization
heuristic. No fixed fractional reduction in computational cost is asserted.

\begin{lemma}[Exact Schmidt rank and single-cut truncation]
\label{lem:schmidt}
For a normalized bipartite vector
$\psi=\sum_{j=1}^r s_j u_j\otimes v_j$ with orthonormal Schmidt vectors,
$s_j>0$, and $\sum_j s_j^2=1$, any MPS bond representing this cut exactly
has dimension at least $r$. Its entropy
$S=-\sum_j s_j^2\log(s_j^2)$ satisfies $S\le\log r$.
If a nonempty retained set has discarded weight
$w=\sum_{j\ \mathrm{discarded}}s_j^2<1$ and $\phi$ is the normalized
retained vector, then
\[
 \norm{\psi-\phi}^2=2(1-\sqrt{1-w})\le2w.
\]
For any Hermitian $H$ on the full bipartite space,
\[
 |\ip{\psi}{H\psi}-\ip{\phi}{H\phi}|
 \le2\norm H\norm{\psi-\phi}\le2\norm H\sqrt{2w}.
\]
\end{lemma}
\begin{proof}
At the cut, a bond of size $D$ factors the coefficient matrix through
$\C^D$, so its matrix rank is at most $D$. The Schmidt decomposition is a
singular-value decomposition with $r$ nonzero singular values, giving
$D\ge r$. Put $p_j=s_j^2$. Concavity of the logarithm gives
$\sum_j p_j\log(1/p_j)\le\log(\sum_j p_j/p_j)=\log r$.
Normalization of the retained part gives
$\ip{\psi}{\phi}=\sqrt{1-w}$, proving the norm identity. Since
$\sqrt{1-w}\ge1-w$ on $[0,1]$, the norm bound follows. Finally,
\[
 \ip{\psi}{H\psi}-\ip{\phi}{H\phi}
 =\ip{\psi-\phi}{H\psi}+\ip{\phi}{H(\psi-\phi)}.
\]
Cauchy--Schwarz and normalization bound the two terms by
$\norm H\norm{\psi-\phi}$ each.
\end{proof}

An entropy upper bound alone does not give an upper bound on exact Schmidt
rank: for any finite $r$, the probabilities
$1-\varepsilon,\varepsilon/(r-1),\ldots,\varepsilon/(r-1)$ have rank
$r$ for $0<\varepsilon<1$ but entropy tending to zero as
$\varepsilon$ tends to zero. The truncation lemma requires a bound on the
actual discarded weight. Repeated truncations require accumulated error
control, and none of these state approximation statements supplies the
complement energy bound in Theorem~\ref{thm:block}.

\subsection{Topology distinguishes connectivity from occupancy}
Suppose a finite set of proposed adsorption positions is supplied and a
threshold $r>0$ defines a simple conflict graph $G_r$: vertices are
positions and an edge means the two positions cannot be occupied together.
For example, an edge can be assigned when their Euclidean distance is at most
$r$. In a model with only these pairwise exclusions, the largest number of
simultaneously occupied vertices is the independence number $\alpha(G_r)$.
The zeroth Betti number $\beta_0$ of the associated Vietoris--Rips complex
counts connected components of its one-skeleton $G_r$.

\begin{proposition}[Component count is generally not occupancy capacity]
\label{prop:graph}
For every finite simple graph $G$, with connected components
$C_1,\ldots,C_m$,
\[
 \alpha(G)=\sum_{j=1}^m\alpha(C_j)\ge m=\beta_0(G).
\]
Equality holds exactly when every component is a complete graph.
\end{proposition}
\begin{proof}
The intersection of an independent set with each component is independent,
so its size is at most $\sum_j\alpha(C_j)$. Conversely, the union of
maximum independent sets from distinct components is independent, because
no edge joins different components. This proves equality in the first
formula. Each nonempty component has at least one vertex, so its independence
number is at least one. A complete component has independence number one.
If a component is not complete, it has two nonadjacent vertices, which form
an independent set of size two. Therefore the sum equals $m$ if and only
if all components are complete. The empty graph satisfies the statement
with both sums zero and the component condition vacuous.
\end{proof}

\begin{example}[A three-position counterexample]
Take positions $(0,0,0)$, $(3,0,0)$, and $(6,0,0)$ in one common length
unit and threshold $r=3.5$. The only edges connect the first to the second
and the second to the third. This three-vertex path is connected, so
$\beta_0=1$. Its two endpoints are nonadjacent and can be occupied together,
whereas all three cannot; hence $\alpha=2$. The example is unchanged if the
edge threshold uses strict instead of non-strict inequality.
\end{example}

Persistent homology may describe connectivity as a distance threshold varies.
It does not prove a capacity rule without an additional physical model. Real
adsorption may also depend on chemical identity, ligand orientation, coverage,
and interactions beyond pairwise exclusion. Even an exact independent set
calculation addresses only the supplied graph. A graph or point cloud whose
reference capacity was assigned by hand is not an independent DFT benchmark,
and a greedy maximal independent set need not be a maximum independent set.
For the path above, choosing the middle vertex first gives a maximal set of
size one although the maximum has size two.

\section{An exactly solvable pairing check}
\label{sec:algebra}
\subsection{Hamiltonian and energy origin}
Reduced pairing models provide a useful algebraic test case in the
Richardson--Gaudin literature \citep{Richardson1963,Dukelsky2004}.
The following two-level calculation is derived directly and needs no general
integrability theorem. It is a toy Hamiltonian, not an ab initio model of
$\mathrm H_2$. No molecular integrals or nuclear repulsion are specified.

Order the four fermionic modes as
$(1\uparrow,1\downarrow,2\uparrow,2\downarrow)$ and restrict to two
electrons. Put
\[
 P_p^\dagger=a_{p\uparrow}^\dagger a_{p\downarrow}^\dagger,
 \quad P_p=a_{p\downarrow}a_{p\uparrow},\quad
 n_{p\sigma}=a_{p\sigma}^\dagger a_{p\sigma},\quad p=1,2.
\]
For real $\varepsilon_1,\varepsilon_2$ and attractive coupling $g\ge0$,
define
\begin{align}
 H_z&=\sum_{p=1}^2\varepsilon_p
       (n_{p\uparrow}+n_{p\downarrow}-1)
       -g\sum_{p,q=1}^2P_p^\dagger P_q,
       \label{eq:pair-z}\\
 H_N&=\sum_{p=1}^2\varepsilon_p
       (n_{p\uparrow}+n_{p\downarrow})
       -g\sum_{p,q=1}^2P_p^\dagger P_q
     =H_z+(\varepsilon_1+\varepsilon_2)I.
       \label{eq:pair-n}
\end{align}
The two conventions differ only by an energy origin, but they must not be
mixed in a secular equation or a numerical comparison.

\begin{proposition}[Complete spectrum in the two-electron sector]
\label{prop:pairing}
Let $d=\varepsilon_1-\varepsilon_2$ and
$\rho=\sqrt{d^2+g^2}$. In the paired basis
$|1\overline1\rangle=P_1^\dagger|0\rangle$,
$|2\overline2\rangle=P_2^\dagger|0\rangle$, the matrix of $H_z$ is
\begin{equation}
 H_z^{\mathrm{pair}}=
 \begin{pmatrix}
 \varepsilon_1-\varepsilon_2-g&-g\\
 -g&-\varepsilon_1+\varepsilon_2-g
 \end{pmatrix}.
 \label{eq:pair-block}
\end{equation}
The complete six-dimensional sector has spectrum, with multiplicity,
\begin{equation}
 \{-g-\rho,\ -g+\rho,\ 0,\ 0,\ 0,\ 0\}.
 \label{eq:pair-spectrum}
\end{equation}
In particular its lowest energy is $E_z^-=-g-\rho$.
The paired energies of $H_N$ are
\begin{equation}
 E_N^\pm=\varepsilon_1+\varepsilon_2-g\pm\rho,
 \label{eq:number-roots}
\end{equation}
and its four unpaired energies are $\varepsilon_1+\varepsilon_2$.
\end{proposition}
\begin{proof}
There are $\binom42=6$ two-electron determinants. Two occupy both spins
of a single orbital; the remaining four place one electron in each orbital.
On $|p\overline p\rangle$, $P_q$ is zero unless $q=p$, in which case
it gives the vacuum with positive sign. This follows directly from the CAR:
$a_{p\downarrow}a_{p\uparrow}a_{p\uparrow}^\dagger
 a_{p\downarrow}^\dagger|0\rangle=|0\rangle$.
Subsequent application of $P_s^\dagger$ gives
$|s\overline s\rangle$. Thus $\sum_{s,q}P_s^\dagger P_q$ has every
entry equal to one in the paired basis. In particular the terms $s=q$
contribute $-g$ to each paired diagonal; they cannot be omitted.
The one-body term in $H_z$ is $d$ on $|1\overline1\rangle$ and
$-d$ on $|2\overline2\rangle$. This proves \eqref{eq:pair-block}.

Every pair annihilator kills an unpaired determinant. Its occupations also
give $n_{p\uparrow}+n_{p\downarrow}-1=0$ for each $p$, so its energy
is zero and it is decoupled from the paired space. Finally,
\[
 \det(H_z^{\mathrm{pair}}-\lambda I)
 =(d-g-\lambda)(-d-g-\lambda)-g^2
 =(\lambda+g)^2-d^2-g^2.
\]
Its roots are $-g\pm\rho$. Since $\rho\ge g\ge0$, the lower root
is at most zero and the upper root is at least zero. Comparing with the four
unpaired zeros establishes the ground energy. Adding the scalar in
\eqref{eq:pair-n} shifts every eigenvalue as stated.
\end{proof}

The upper paired root is not generally the first excited level of the full
sector: the unpaired zeros must also be included. When $g>0$, a normalized
ground vector in the paired space is
\[
 c^- =\frac{(g,d+\rho)^T}{\sqrt{g^2+(d+\rho)^2}}.
\]
Indeed the first row of
$(H_z^{\mathrm{pair}}-E_z^-I)(g,d+\rho)^T$ vanishes by cancellation,
and the second is $-g^2+(\rho-d)(\rho+d)=0$.
The denominator is positive for $g>0$. Its orthogonal vector
$c^+=(-c^-_2,c^-_1)^T$ is the upper eigenvector: symmetry makes the
orthogonal one-dimensional subspace invariant, and the trace fixes its
eigenvalue to $-g+\rho$. For $g=0$, determinant basis vectors diagonalize
the block. If also $d=0$, both paired energies and all unpaired energies
coincide at zero. These cases should be handled directly rather than by
dividing by $g$.

\subsection{The one-pair secular equation and its excluded points}
In the number convention,
$H_N^{\mathrm{pair}}=\operatorname{diag}(2\varepsilon_1,
2\varepsilon_2)-g\boldsymbol 1\boldsymbol 1^T$.
An eigenvector $(c_1,c_2)^T$ with energy $E$ obeys
\[
 (2\varepsilon_p-E)c_p=g(c_1+c_2),\qquad p=1,2.
\]
If neither denominator vanishes, a nonzero vector must have
$c_1+c_2\ne0$; otherwise both coefficients would be zero. Dividing and
summing therefore gives
\begin{equation}
 1=g\left(\frac{1}{2\varepsilon_1-E}
          +\frac{1}{2\varepsilon_2-E}\right).
 \label{eq:richardson}
\end{equation}
Conversely, a finite $E$ satisfying this equation away from its poles yields
an eigenvector with coefficients proportional to
$(2\varepsilon_p-E)^{-1}$, by substitution in the two linear equations.
The polynomial form is
\begin{equation}
 (2\varepsilon_1-E)(2\varepsilon_2-E)
 -g\bigl[(2\varepsilon_1-E)+(2\varepsilon_2-E)\bigr]=0.
 \label{eq:pair-polynomial}
\end{equation}
This is the determinant equation and remains meaningful where the rational
form is undefined. It produces exactly \eqref{eq:number-roots}.
At $g=0$ the eigenvalues lie at poles of the rational expression. At
$\varepsilon_1=\varepsilon_2=\varepsilon$ and $g>0$, the upper state
$(1,-1)^T/\sqrt2$ has $E=2\varepsilon$ and zero coefficient sum;
it too is excluded from the rational derivation. A pole is not a zero
residual. Near a pole, the polynomial or matrix eigenvector residual is the
appropriate numerical check, with tolerances scaled to the matrix norm.

\subsection{Relation to the energy certificate}
The paired block is itself a scalar active/complement split with
$a=d-g$, $c=-d-g$, and $b=g$. Its smaller eigenvalue equals the bound
in Theorem~\ref{thm:block}, providing a sharp algebraic check. If the lower
diagonal entry is chosen as the active block and $d\ne0$, the complementary
diagonal lies higher and Corollary~\ref{cor:gap} also applies to this paired
problem. To claim the ground of the full sector one must still account for
the unpaired block; Proposition~\ref{prop:pairing} does so explicitly.

Agreement between a CAR matrix construction, the roots above, and an
independent Jordan--Wigner construction tests this toy algebra. It supplies no
error estimate between a molecular Coulomb Hamiltonian and a reduced pairing
Hamiltonian. Establishing such a reduction would require specified molecular
integrals and a certified bound on the omitted interactions. General claims
about Yangian symmetry, Howe duality, or exact molecular solvability are not
needed for this example and are not established here.

\section{From electronic intervals to conditional kinetic bounds}
\label{sec:approximations}
\subsection{Barrier intervals and an explicit error budget}
Let $R$ and $\mathrm{TS}$ denote two specified geometries, with consistent
stoichiometry, charge convention, and energy reference. The label
$\mathrm{TS}$ is only a label until a saddle point and its connection to the
reactant basin have been established. Suppose
\[
 L_R\le E_R\le U_R,\qquad
 L_{\mathrm{TS}}\le E_{\mathrm{TS}}\le U_{\mathrm{TS}}.
\]

\begin{proposition}[Interval subtraction]
\label{prop:barrier}
The electronic energy difference $B_{\mathrm{el}}=E_{\mathrm{TS}}-E_R$
satisfies
\begin{equation}
 L_{\mathrm{TS}}-U_R\le B_{\mathrm{el}}
 \le U_{\mathrm{TS}}-L_R.
 \label{eq:barrier}
\end{equation}
These endpoints are optimal given only the two independent intervals.
If an additive correction $\Delta$ obeys
$\Delta\in[\Delta_L,\Delta_U]$, then
\begin{equation}
 B=B_{\mathrm{el}}+\Delta\in
 [L_{\mathrm{TS}}-U_R+\Delta_L,
  U_{\mathrm{TS}}-L_R+\Delta_U].
 \label{eq:corrected-barrier}
\end{equation}
\end{proposition}
\begin{proof}
The inequalities $E_{\mathrm{TS}}\ge L_{\mathrm{TS}}$ and
$-E_R\ge-U_R$ give the lower endpoint on addition. Adding
$E_{\mathrm{TS}}\le U_{\mathrm{TS}}$ and $-E_R\le-L_R$ gives the
upper endpoint. The choices $(E_{\mathrm{TS}},E_R)=
(L_{\mathrm{TS}},U_R)$ and $(U_{\mathrm{TS}},L_R)$ attain the respective
endpoints, proving optimality without further joint information. Adding the
bounds for $\Delta$ proves \eqref{eq:corrected-barrier}.
\end{proof}

The interval width in \eqref{eq:barrier} is
$(U_{\mathrm{TS}}-L_{\mathrm{TS}})+(U_R-L_R)$.
Correlated energy errors can sometimes yield a narrower difference bound,
but their correlation must be proved. Cancellation cannot simply be assumed.
For a free-energy barrier, $\Delta$ must include the relevant zero-point,
thermal, entropic, environmental, and standard-state corrections with their
own justified bounds. A negative lower endpoint is allowed algebraically;
clipping it to zero requires a separate proof that the target barrier is
nonnegative.

\begin{lemma}[Finite-model perturbation]
\label{lem:perturbation}
If $H$ and $\widetilde H$ are Hermitian on the same finite space and
$\norm{H-\widetilde H}\le\delta$, then
$|E_0(H)-E_0(\widetilde H)|\le\delta$. Consequently, an interval
$[L,U]$ for $E_0(\widetilde H)$ yields $[L-\delta,U+\delta]$ for
$E_0(H)$.
\end{lemma}
\begin{proof}
Every normalized $v$ satisfies
$|\ip{v}{(H-\widetilde H)v}|\le\delta$. Inserting a lowest eigenvector
of $\widetilde H$ into the variational principle for $H$ gives
$E_0(H)\le E_0(\widetilde H)+\delta$. Interchanging the two matrices
gives the reverse bound. Enlarging each interval endpoint by $\delta$
then contains the target energy.
\end{proof}

This lemma can absorb rigorously bounded integral or model perturbations
only after the spaces and energy conventions have been matched. It does not
supply $\delta$ automatically. For a conforming finite-basis restriction of
a continuum Hamiltonian, the variational principle gives
$E_{\mathrm{continuum}}\le E_{\mathrm{basis}}$. A lower bound on
$E_{\mathrm{basis}}$ therefore need not be a lower bound on
$E_{\mathrm{continuum}}$. Basis-set, slab-size, embedding, and boundary
errors each require a separate comparison theorem or error estimate.

\subsection{Geometry, paths, and Born--Oppenheimer assumptions}
Computing an energy at a candidate saddle does not bound the error in its
location. If an energy surface $E(q)$ is known to be Lipschitz with constant
$M$ on a domain and a target geometry $q_*$ is known to satisfy
$\norm{q_*-q_0}\le r$, then $|E(q_*)-E(q_0)|\le Mr$ by the
Lipschitz definition. Both $M$ and the location bound are additional inputs.
A small gradient at $q_0$ alone does not establish a particular saddle or
its connection to a reactant.

\begin{proposition}[What path samples can certify]
\label{prop:path}
Let $E:[0,1]\to\R$ be $M$-Lipschitz. Suppose samples $s_i$ form a
covering with radius $h$, meaning that every $s\in[0,1]$ has
$|s-s_i|\le h$ for some $i$, and certified sample intervals satisfy
$L_i\le E(s_i)\le U_i$. Then
\begin{equation}
 \max_i L_i\le\max_{s\in[0,1]} E(s)\le\max_i U_i+Mh.
 \label{eq:path}
\end{equation}
Without a modulus of continuity or other control between samples, finitely
many sampled values do not give any finite uniform upper bound.
\end{proposition}
\begin{proof}
Continuity and compactness ensure the maximum exists. Each sampled value is
at most that maximum, proving the lower bound. For any $s$, choose a
sample within $h$. Then
$E(s)\le E(s_i)+M|s-s_i|\le\max_i U_i+Mh$; maximizing gives the
upper bound. For the final statement, choose a nonempty open interval
avoiding all of the finitely many sample points. A continuous triangular
bump supported in that interval vanishes at every sample and can have an
arbitrarily large positive height. Adding it preserves all sampled values
while making the maximum arbitrarily large.
\end{proof}

Even a certified maximum along one path is not automatically the minimum
barrier over all possible paths, and endpoint energies alone do not certify
a path maximum. Likewise, finitely sampled electronic gaps do not establish
a uniform positive gap between samples.

Rigorous Born--Oppenheimer results require their own setting and hypotheses;
for example, Hagedorn's analysis treats high-order corrections for diatomic
Coulomb bound states \citep{Hagedorn1988}. That result is not invoked here
as a numerical surface-reaction error bound. A mass ratio such as
$\varepsilon=(m_e/M)^{1/4}$ is a dimensionless parameter. Even if an
applicable theorem gives $|\delta E|\le C\varepsilon^p$ on a specified
domain, one must bound $C$, specify energy units and the approximation order,
and check its spectral and regularity assumptions. Evaluating
$\varepsilon^5$ and multiplying by an assumed reaction-energy scale does
not perform these steps.

The role of electronic derivatives can be seen without any asymptotic
approximation. For a differentiable Hermitian family $H_e(q)$ with
differentiable normalized eigenvectors $\phi_m(q)$ and distinct energies,
differentiation of $H_e\phi_m=E_m\phi_m$ gives
\[
 (\partial_q H_e)\phi_m+H_e\partial_q\phi_m
 =(\partial_qE_m)\phi_m+E_m\partial_q\phi_m.
\]
Taking the inner product with $\phi_n$, $n\ne m$, and using
orthogonality yields
\begin{equation}
 \ip{\phi_n}{\partial_q\phi_m}
 =\frac{\ip{\phi_n}{(\partial_qH_e)\phi_m}}{E_m-E_n}.
 \label{eq:derivative-coupling}
\end{equation}
Thus a gap alone does not determine derivative coupling; the numerator and
its units matter. In particular, an orbital HOMO--LUMO gap at one geometry
is not a certified many-electron gap along a reaction path. No absence of
crossings, uniform nonadiabatic error, or CaF2 surface-reaction BO bound is
proved by this manuscript.

\subsection{Conditional transition-state rate intervals}
The connection to kinetics is an assumed transition-state or activated-rate
model motivated by transition-state theory \citep{Eyring1935}:
\begin{equation}
 k=A\exp(-B/\theta),\qquad \theta>0.
 \label{eq:tst}
\end{equation}
Here $\theta$ has the same energy units as $B$: it may be $k_{\mathrm B}T$
for energies per event or $RT$ for molar energies. The prefactor must have
the rate units appropriate to the selected kinetic convention. For a
unimolecular Eyring model it may include $k_{\mathrm B}T/h$ and a
transmission coefficient; other conventions can require concentration or
standard-state factors. These are model inputs, not consequences of the
electronic theorem. The results below assume a fixed known $\theta$.

\begin{theorem}[Monotone rate propagation]
\label{thm:rate}
Assume \eqref{eq:tst}, $B\in[L,U]$, and
$0<A_L\le A\le A_U$. Then
\begin{equation}
 0<k_L:=A_L e^{-U/\theta}\le k\le A_U e^{-L/\theta}=:k_U.
 \label{eq:rate-interval}
\end{equation}
The endpoints are the extremal rates over the rectangular uncertainty set
$[A_L,A_U]\times[L,U]$.
\end{theorem}
\begin{proof}
Since $\theta>0$, the barrier inequalities give
$-U/\theta\le-B/\theta\le-L/\theta$.
The exponential is positive and increasing, so
$e^{-U/\theta}\le e^{-B/\theta}\le e^{-L/\theta}$.
Multiplication by the positive prefactor bounds proves
\eqref{eq:rate-interval}. The parameter choices $(A_L,U)$ and $(A_U,L)$
attain the lower and upper endpoints.
\end{proof}

No independence assumption on random variables is involved: these are
worst-case deterministic bounds. If the prefactor and barrier are known to
obey an additional joint constraint, the rectangular bounds remain sound
but may cease to be sharp. A change in $\theta$ requires its own interval
analysis; it cannot be absorbed silently into a fixed-temperature result.

\begin{corollary}[Robust comparison of two channels]
\label{cor:ranking}
For two rates at the same $\theta$, suppose the assumed model gives a common
positive prefactor $A$. If their barriers satisfy $U_1<L_2$, then
$k_1>k_2$ for every permitted barrier pair. With potentially different
positive prefactor intervals, the sufficient condition is
\begin{equation}
 A_{1,L}e^{-U_1/\theta}>A_{2,U}e^{-L_2/\theta},
 \quad\text{equivalently}\quad
 L_2-U_1>\theta\log\!\left(\frac{A_{2,U}}{A_{1,L}}\right).
 \label{eq:robust-ranking}
\end{equation}
\end{corollary}
\begin{proof}
For the common prefactor, $B_1\le U_1<L_2\le B_2$ and the strict
monotonic decrease of $e^{-B/\theta}$ give $k_1>k_2$.
In the general case Theorem~\ref{thm:rate} gives
$k_1\ge A_{1,L}e^{-U_1/\theta}$ and
$k_2\le A_{2,U}e^{-L_2/\theta}$; the displayed strict separation
therefore suffices. Taking logarithms of its positive sides and rearranging
proves the equivalent condition.
\end{proof}

Touching intervals give at most a non-strict conclusion. Overlapping rate
intervals leave the ordering unresolved by this certificate; they do not
show that the true rates coincide. Ranking elementary rates is also distinct
from ranking net reaction fluxes, which depend on populations and the network.

\begin{proposition}[Energy precision is set by logarithmic rate uncertainty]
\label{prop:log-budget}
For the endpoints in \eqref{eq:rate-interval},
\begin{equation}
 \log(k_U/k_L)=\log(A_U/A_L)+(U-L)/\theta.
 \label{eq:log-width}
\end{equation}
Consequently the interval has ratio $k_U/k_L\le F$, for $F\ge1$, if and
only if
$\log(A_U/A_L)+(U-L)/\theta\le\log F$.
Alternatively, if $\widehat k=\widehat A e^{-\widehat B/\theta}$,
$|\log(A/\widehat A)|\le\alpha$, and
$|B-\widehat B|\le\delta$, then
\begin{equation}
 |\log(k/\widehat k)|\le\alpha+\delta/\theta.
 \label{eq:log-error}
\end{equation}
\end{proposition}
\begin{proof}
Dividing the positive endpoints gives
$k_U/k_L=(A_U/A_L)e^{(U-L)/\theta}$; taking the logarithm proves the
identity and the equivalence. Similarly,
$\log(k/\widehat k)=\log(A/\widehat A)-(B-\widehat B)/\theta$,
and the triangle inequality gives \eqref{eq:log-error}.
\end{proof}

To guarantee a multiplicative factor $F>1$ relative to $\widehat k$, a
sufficient budget is $\alpha+\delta/\theta\le\log F$.
If $\alpha>\log F$, energy refinement alone cannot satisfy this budget.
For a relative error tolerance $0<\epsilon<1$, the stronger condition
$|\log(k/\widehat k)|\le\log(1+\epsilon)$ implies
$|k/\widehat k-1|\le\epsilon$: the ratio lies between
$1/(1+\epsilon)$ and $1+\epsilon$, and
$1-1/(1+\epsilon)=\epsilon/(1+\epsilon)\le\epsilon$.
These formulas provide a target-dependent energy tolerance instead of a
universal claim of chemical accuracy.

\subsection{Finite-state kinetic stability with a complete stated generator}
Let $S=\{1,\ldots,m\}$ be a nonempty finite state space. A row-generator
$Q$ satisfies $Q_{ij}\ge0$ for $i\ne j$ and
$Q_{ii}=-\sum_{j\ne i}Q_{ij}$. Row distributions evolve as
$p(t)=p(0)e^{tQ}$. The following result compares two fully specified,
time-homogeneous generators on the same state space. It makes no
irreducibility or reversibility assumption.

\begin{lemma}[Stochastic semigroups and contraction]
\label{lem:markov}
For such a generator, $P_t=e^{tQ}$ has nonnegative entries and every row
sums to one for $t\ge0$. For any real row vector $z$,
$\norm{zP_t}_1\le\norm z_1$.
\end{lemma}
\begin{proof}
Choose a real $\nu>0$ with $\nu\ge\max_i(-Q_{ii})$ and put
$P=I+Q/\nu$. Off-diagonal entries and diagonal entries of $P$ are
nonnegative and its rows sum to one. Products of such matrices have the
same properties. Since $Q=\nu(P-I)$ and $I$ commutes with $P$,
\[
 e^{tQ}=e^{-\nu t}\sum_{n=0}^\infty\frac{(\nu t)^n}{n!}P^n.
\]
The coefficients are nonnegative and sum to one. Entrywise convergence
therefore proves the stochastic property. Finally,
\[
 \sum_j\left|\sum_i z_i(P_t)_{ij}\right|
 \le\sum_i|z_i|\sum_j(P_t)_{ij}=\sum_i|z_i|,
\]
which is the asserted contraction.
\end{proof}

\begin{theorem}[Finite-state generator perturbation bound]
\label{thm:ctmc}
Let $Q,\widehat Q$ be row-generators on $S$ and define
\begin{equation}
 \eta=\frac12\max_i\sum_{j=1}^m|Q_{ij}-\widehat Q_{ij}|.
 \label{eq:eta}
\end{equation}
The sum includes the diagonal entry. For the same initial probability
row vector $p$, put $p_t=pe^{tQ}$ and
$\widehat p_t=pe^{t\widehat Q}$. Then, for every $t\ge0$,
\begin{equation}
 \TV(p_t,\widehat p_t):=\tfrac12\norm{p_t-\widehat p_t}_1
 \le\min(1,t\eta).
 \label{eq:ctmc-tv}
\end{equation}
For any real observable $f:S\to\R$, define
$\operatorname{range}(f)=\max_i f(i)-\min_i f(i)$. Then
\begin{equation}
 |p_t f-\widehat p_t f|
 \le\operatorname{range}(f)\min(1,t\eta).
 \label{eq:observable}
\end{equation}
If the initial distributions differ, the corresponding bound is
$\min(1,\TV(p,\widehat p)+t\eta)$.
\end{theorem}
\begin{proof}
All matrix exponentials are differentiable in finite dimension. For a fixed
$t$, differentiate $F(s)=e^{sQ}e^{(t-s)\widehat Q}$ on $[0,t]$.
Commutation of a matrix with its own exponential gives
\[
 F'(s)=e^{sQ}(Q-\widehat Q)e^{(t-s)\widehat Q}.
\]
Integrating from $0$ to $t$ yields the Duhamel identity
\begin{equation}
 e^{tQ}-e^{t\widehat Q}
 =\int_0^t e^{sQ}(Q-\widehat Q)e^{(t-s)\widehat Q}\,ds.
 \label{eq:duhamel}
\end{equation}
For every probability vector $v$, the triangle inequality gives
\[
 \norm{v(Q-\widehat Q)}_1
 \le\sum_i v_i\sum_j|Q_{ij}-\widehat Q_{ij}|\le2\eta.
\]
By Lemma~\ref{lem:markov}, $pe^{sQ}$ is a probability vector, and
right multiplication by $e^{(t-s)\widehat Q}$ is an $\ell^1$
contraction. Left-multiplying \eqref{eq:duhamel} by $p$ and applying
these two estimates under the integral gives
$\norm{p_t-\widehat p_t}_1\le\int_0^t2\eta\,ds=2t\eta$.
Two probability vectors have $\ell^1$ distance at most two, proving
\eqref{eq:ctmc-tv}.

For the observable inequality, put $z=p_t-\widehat p_t$ and
$c=(\max_i f(i)+\min_i f(i))/2$. Since $\sum_i z_i=0$,
\[
 |zf|=\left|\sum_i z_i(f(i)-c)\right|
 \le\frac{\operatorname{range}(f)}2\sum_i|z_i|
 =\operatorname{range}(f)\TV(p_t,\widehat p_t).
\]
Combining with \eqref{eq:ctmc-tv} proves \eqref{eq:observable}.
For different initial distributions, write
\[
 pe^{tQ}-\widehat p e^{t\widehat Q}
 =p(e^{tQ}-e^{t\widehat Q})
   +(p-\widehat p)e^{t\widehat Q}.
\]
The first term has norm at most $2t\eta$, and contraction bounds the
second by $\norm{p-\widehat p}_1$. Division by two and the universal
probability-distance bound complete the proof.
\end{proof}

\begin{corollary}[From individual rate intervals to a generator bound]
\label{cor:rates-generator}
Suppose all off-diagonal entries are specified and have validated intervals
$Q_{ij}\in[l_{ij},u_{ij}]$ with $0\le l_{ij}\le u_{ij}$.
Choose $\widehat Q_{ij}\in[l_{ij},u_{ij}]$ and set both diagonals by
zero row sums. Define
$\epsilon_{ij}=\max(\widehat Q_{ij}-l_{ij},u_{ij}-\widehat Q_{ij})$
for $i\ne j$. Then
\begin{equation}
 \eta\le\max_i\sum_{j\ne i}\epsilon_{ij}.
 \label{eq:rate-generator}
\end{equation}
\end{corollary}
\begin{proof}
Each off-diagonal difference is at most $\epsilon_{ij}$ in absolute
value, and
$|Q_{ii}-\widehat Q_{ii}|=
|\sum_{j\ne i}(Q_{ij}-\widehat Q_{ij})|
\le\sum_{j\ne i}\epsilon_{ij}$.
The full row sum of absolute differences is therefore at most twice this
sum. Substitution in \eqref{eq:eta} proves the claim.
\end{proof}

An entry may be the sum of several event rates; its interval must include
all those contributions. A nominally absent transition is justified only if
its rate is zero by the stated model or is bounded within the error budget.
The theorem does not certify that an event catalog contains every relevant
physical reaction. It also does not compare generators on different state
spaces without an additional embedding or coarse-graining argument. Pulsed
operating conditions and transport can require time-dependent generators or
a larger model, which must be analyzed separately.

A bounded terminal occupancy is an observable covered by
\eqref{eq:observable}. Growth per cycle is not automatically such an
observable: it may involve accumulated event counts, changing surface states,
transport, and a conversion from deposited material to thickness. A pathwise
count is not generally bounded by the range of a state function on $S$.
Thus the CTMC result is a conditional finite-state stability theorem, not a
direct growth-per-cycle guarantee. Establishing the physical network and the
observable mapping remains a separate task.

\section{Computational scope and unresolved tasks}
\label{sec:complexity}
\subsection{An enclosure theorem is not a complexity breakthrough}
The finite-basis electronic-structure problem in the setting of
\citet{OGorman2022} is a worst-case QMA-complete ground-energy promise
problem. The relevant precision is an inverse-polynomial additive promise
gap under the paper's input and scaling conventions, not a universal relative
error guarantee for all molecular energies. Fermionic $N$-representability
has a corresponding QMA-completeness result \citep{Liu2007}.
These classifications concern families of encoded instances. They neither
prove that every molecule requires exponential time nor provide a probability
measure under which most molecules are easy. No separation of complexity
classes is assumed as a proved theorem here.

The block estimate reduces the \emph{verification statement} to a few scalar
inequalities once their hypotheses are supplied. It does not reduce the cost
of constructing those hypotheses. The active eigenproblem, the complementary
lower bound, the coupling norm, integral generation, and verified rounding all
contribute to the cost. An active-space dimension that is small relative to
the full determinant space need not be small as a function of the input size.
Moreover, a complement estimate may be cheap and so loose that the final
barrier interval is inconclusive.

The same issue arises for RDM and tensor representations. The dimensions of
an explicit $p$-RDM on antisymmetric $p$-tuples are
$\binom Kp\times\binom Kp$. For fixed $p$, this gives a polynomial
number of entries, but does not make exact $N$-representability easy. For
commuting SOS variables, if the number of variables is $m$, the number of
monomials of total degree at most $d$ is $\binom{m+d}{d}$: append a slack
exponent to the $m$ nonnegative exponents so their sum is exactly $d$, and
count weak compositions into $m+1$ parts. In particular, taking $m=K^4$
and $d=3$ gives
\[
 \binom{K^4+3}{3}
 =\frac{(K^4+3)(K^4+2)(K^4+1)}6=\Theta(K^{12}),
\]
not $O(K^6)$. This count is only for that explicitly chosen commuting
monomial basis; a fermionic operator construction must define its own basis.
A dimension count alone supplies neither an SDP runtime nor a convergence
rate, conditioning estimate, or accuracy guarantee.

\subsection{What an average-case statement would require}
An average-case complexity claim needs a family of probability distributions
$\mathcal D_n$ on finite encoded inputs, an algorithm, a success criterion,
an accuracy target, and a cost model. Qualitative labels such as stable,
local, or chemically typical do not specify a normalized distribution.
Empirical success on a selected set is evidence about that set, not a proof
of polynomial expected runtime.

\begin{proposition}[A sufficient tail condition, not an established chemical law]
\label{prop:average}
Let $T_n\ge0$ be the runtime of a specified algorithm under a specified
$\mathcal D_n$, where $n\ge1$. Suppose constants $C,C',c,d>0$ and
$k\ge0$ satisfy: outside an event of probability at most $C'e^{-cn}$,
$T_n\le Cn^k$, and on every input $T_n\le Ce^{dn}$. Then
\[
 \mathbb E[T_n]\le Cn^k+CC'e^{(d-c)n}.
\]
In particular these hypotheses imply a polynomial expectation if $c\ge d$.
\end{proposition}
\begin{proof}
Let $A_n$ be the exceptional event. Split the expectation as
$\mathbb E[T_n\boldsymbol1_{A_n^c}]+
\mathbb E[T_n\boldsymbol1_{A_n}]$.
The first term is at most $Cn^k\mathbb P(A_n^c)\le Cn^k$, and the
second is at most $Ce^{dn}\mathbb P(A_n)\le CC'e^{(d-c)n}$.
If $c\ge d$, the second bound is at most $CC'$.
\end{proof}

An exponentially small difficult subset is therefore not by itself a
polynomial expectation theorem: the runtime on it also matters. For example,
a runtime equal to $1$ outside an event of probability $e^{-cn}$ and to
$e^{dn}$ on that event has expectation
$1-e^{-cn}+e^{(d-c)n}$, which grows exponentially when $d>c$.
No distribution or tail estimate meeting these hypotheses has been
constructed here for surface chemistry. No percentages of easy molecules or
timelines for overcoming an exponential cost are inferred.

\subsection{A research protocol for a specified elementary step}
The results suggest a concrete sequence of tasks with explicit stopping
conditions. The aim is an informative interval, which may be too wide to
rank pathways; a failed ranking is not an algorithmic or physical prediction.

\begin{enumerate}
\item \textbf{Specify the physical question.}
For $\mathrm{Al_2O_3}/\mathrm{MoCl_5}$, record the substrate phase and
facet, termination and hydroxylation, slab size and boundary conditions,
coverage, precursor charge and spin, co-reactant, temperature, and reference
reservoirs. For CaF2 crystal deposition, resolve the precursor pair, the
co-reactant, the substrate phase and facet, and the reaction stoichiometry
before defining any elementary step.
Define which elementary steps and competing channels are to be compared.

\item \textbf{Fix a reproducible finite electronic target.}
Provide geometries, orbital basis and ordering, electron count, symmetry
sector, Hamiltonian coefficients, nuclear terms, and units. Record how each
geometry was obtained and what establishes a reactant minimum or a connected
transition state. A Hamiltonian at an unverified candidate geometry supports
only a statement about that geometry.

\item \textbf{Produce independently justified endpoints.}
Construct a normalized trial state and a bound on its expectation. Supply
$\underline a,\underline c,\overline b$ for the selected block split,
or a dual-feasible outer-relaxation certificate, with all feasibility and
rounding errors accounted for. Report the cost of the complement information
separately. Compare with full diagonalization only where the latter is
feasible, without treating the comparison itself as a scalable method.

\item \textbf{Bound every correction to the requested target.}
Use interval subtraction for electronic differences. Add justified bounds
for geometry, thermodynamic, environmental, and other corrections that are
needed for the chosen barrier definition. If a continuous path maximum is
required, supply between-sample control as in Proposition~\ref{prop:path}.
Do not infer a free-energy or nonadiabatic error bar from electronic solver
convergence alone.

\item \textbf{Set the precision from the kinetic decision.}
State the rate model, temperature, energy-unit convention, and positive
prefactor interval. Use Proposition~\ref{prop:log-budget} to decide whether
the energy precision can meet the desired rate ratio. Apply
Corollary~\ref{cor:ranking} only when its strict separation holds; otherwise
report an unresolved ordering and identify the dominant interval width.

\item \textbf{Define the network and observable separately.}
For a finite-state kinetic model, enumerate states, transitions, environmental
conditions, and all generator entries. Establish a rate bound for every entry
and use Corollary~\ref{cor:rates-generator}. A bounded state observable then
inherits Theorem~\ref{thm:ctmc}. Any relation to accumulated deposition or
film thickness requires its own model and error analysis. Validate physical
inputs against independently sourced data when available.
\end{enumerate}

\subsection{Open mathematical and computational questions}
The central open problem is a useful complementary-space lower bound whose
construction avoids enumerating the full omitted determinant space. Candidate
routes include decompositions into positive operators, valid RDM dual
certificates, or special operator structure. Each route must exhibit a
checkable inequality for the actual complement, together with its cost and
its loss of precision. An upper estimate for that complement cannot substitute
for the required lower bound.

A second problem is to combine active-state improvement and lower-bound
improvement adaptively. Widening an active space may improve a variational
upper bound, but a useful algorithm must show how to maintain and evaluate
the lower certificate as the partition changes. Localized orbitals, symmetry,
or tensor compression can aid the search without proving a uniform rate of
convergence or a polynomial runtime.

A third problem is validated numerical evaluation. Floating-point eigenvalues,
SDP feasibility residuals, and rate exponentials must be enclosed with error
bounds tight enough to preserve the desired energy and rate inequalities.
The residual correction in Proposition~\ref{prop:dual} gives one possible
pattern, but its spectral residual and trace bound still have to be verified
for the concrete SDP.

Finally, uniform geometry and path errors, a defensible thermodynamic
correction budget, a complete physical event catalog, and a measured
observable mapping remain unresolved for the named ALD targets. The present
paper provides conditional implications, not solutions to these input
problems. Its next substantive result should be a fully specified example
where each hypothesis can be checked and where the final interval either
supports or demonstrably fails to support a stated chemical decision.

\section{A worked protocol for one elementary step}
\label{sec:worked}

Section~\ref{sec:intro} listed the inputs a certificate needs. This section
supplies them for one system far enough to expose where the method is already
informative and where it is not. Nothing below is a molecular energy
calculation; the certified content is arithmetic on published growth data plus
two geometric inequalities, and the negative result is proved.

\subsection{System and evidence status}

The target is $\mathrm{CaF_2}$ growth from a calcium amidinate and hydrogen
fluoride on the fluorine-terminated $\mathrm{CaF_2}(111)$ cleavage plane.
Every choice below is either taken from a cited measurement or is explicitly
labelled as a modelling decision.

\begin{itemize}
\item \emph{Surface.} The $(111)$ plane is the natural cleavage plane and is
F-terminated with a $(1\times1)$ unreconstructed termination
\citep{Vogt2005,Puchin2001}. Its repeat unit along $[111]$ is a
F--Ca--F trilayer of height $a_0/\sqrt3$ \citep{Calleja2005}.
\item \emph{Lattice constant.} $a_0 = 5.4626$~\AA, room-temperature XRD powder
data (Swanson \& Tatge, NBS Circular 539 (1953) 69--70; ICSD 52754).
\item \emph{Precursor pair.} Bis($N,N'$-di-\emph{i}-propylformamidinato)calcium(II)
with anhydrous HF, reported at 175--250~$^\circ$C with a growth per cycle of
$0.3$--$0.4$~\AA/cycle \citep{Hennessy2023}.
\item \emph{Precursor nuclearity.} The compound is a \emph{dimer} in the solid
state and is retained in solution, with Ca$\cdots$Ca $=3.323$~\AA, $\mathrm{Ca_2C_{28}H_{60}N_8}$,
$M=589.00$~g\,mol$^{-1}$, monoclinic $P2_1/c$, $V=3754(3)$~\AA$^3$, $Z=4$ at
100~K \citep{Kim2016}. This matters below, because the steric footprint of a
dimer is not that of a monomer.
\item \emph{Adsorbate state.} Periodic DFT places HF \emph{molecularly} on
$\mathrm{CaF_2}(111)$ without dissociation, with $E_{\mathrm{ads}}=-0.40$,
$-0.48$ and $-0.62$~eV at 100\%, 50\% and 25\% coverage, the H hydrogen-bonded
to a surface F \citep{Kaawar2019}. Dissociation to a surface bifluoride
$\mathrm{FHF^-}$ is reported on $(100)$, not $(111)$.
\item \emph{Competing routes in the literature.} $\mathrm{Ca(thd)_2}$ with
$\mathrm{TiF_4}$ at 300--450~$^\circ$C \citep{Pilvi2007,Pilvi2008}; HF--pyridine
as a fluorine co-reactant \citep{Lee2016}; fluorinated $\beta$-diketonates as
single-source precursors \citep{LoPresti2022}; the first metal-fluoride ALD
process \citep{Ylilammi1994}.
\end{itemize}

The general mechanism is better constrained than it first appears, even though
nothing is known specifically for this precursor pair. In situ quartz-crystal
microbalance and infrared work on $\mathrm{AlF_3}$ from $\mathrm{Al(CH_3)_3}$ and
HF finds that the HF pulse both protonates the metal--ligand bond
($\mathrm{HF}+\mathrm{M{-}CH_3}\to\mathrm{M{-}F}+\mathrm{CH_4}$) \emph{and}
leaves molecular HF adsorbed as a surface fluorine reservoir whose coverage
fixes the cycle stoichiometry \citep{Lee2015}. The same scheme is written for
$\mathrm{ZrF_4}$ from a tetrakis(ethylmethylamido)zirconium precursor, where the
volatile byproduct is the free amine \citep{Lee2016}. Protonating an amidinate
likewise gives the conjugate amidine, and amidine desorption has been observed
from a metal surface, together with re-adsorption that is a known carbon
contamination route \citep{Ma2011}.

For the present system, therefore, the expected net reaction is
\[
\mathrm{Ca(amd)_2} + 2\,\mathrm{HF}
  \longrightarrow \mathrm{CaF_2} + 2\,\mathrm{Hamd},
\]
with $\mathrm{Hamd} = N,N'$-di-\emph{i}-propylformamidine. No in situ
mass-balance, infrared, or mass-spectrometric study of this pair exists, and
\emph{no measured or computed activation energy for HF reacting with any metal
amidinate surface has been found}. The only activation energies available for
HF fluorination concern amorphous alumina
($0.18$--$1.14$~eV) \citep{Khumaini2025}; they are not transferable to
$\mathrm{CaF_2}$ and are not used below.

\subsection{Certified arithmetic from published growth data}

Two published numbers, one measured and one geometric, determine the fraction
of a layer deposited per cycle. Let $p$ be the trilayer height and $s$ the
growth per cycle. Writing $s = p\,c$ with $c$ the fractional coverage, an
interval $s\in[s_1,s_2]$ with $p>0$ gives $c\in[s_1/p,\ s_2/p]$
(Lemma~\ref{lem:coverage-lean}).

\begin{lemma}[Interval coverage]
\label{lem:coverage-lean}
For real $p>0$ and $s=p\,c$, if $s_1\le s\le s_2$ then
$s_1/p\le c\le s_2/p$.
\end{lemma}
\begin{proof}
From $s=p\,c$ we get $c=s/p$. Since $p>0$, division by $p$ is monotone, so
$s_1/p\le s/p\le s_2/p$.
\end{proof}

With $p=a_0/\sqrt3=3.1538$~\AA{} and $s\in[0.3,0.4]$~\AA/cycle,
\[
c \in \left[\frac{0.3}{3.1538},\ \frac{0.4}{3.1538}\right]
  = [9.51\%,\ 12.68\%].
\]
Each cycle therefore completes about one tenth of a trilayer.

On the unreconstructed F-terminated $(111)$ plane the Ca ions form a triangular
lattice of spacing $a_0/\sqrt2$, so one Ca site occupies
\[
A_{\mathrm{site}} = \tfrac{\sqrt3}{2}\left(\tfrac{a_0}{\sqrt2}\right)^2
= \tfrac{\sqrt3}{4}a_0^2 = 12.9211\ \text{\AA}^2 .
\]
The areal density of deposited Ca follows as
$\sigma = c/A_{\mathrm{site}} \in [0.00736,\ 0.00982]$~\AA$^{-2}$, i.e.\ one
deposited Ca per $101.9$--$135.8$~\AA$^2$ of surface. That is $7.9$ to $10.5$
times the area of a bare Ca site.

\begin{remark}
This ratio is the quantitative statement that the cycle is \emph{not} limited
by the availability of surface sites. It does not identify the limiting
mechanism.
\end{remark}

\subsection{A geometric design rule for the ligand}

The ratio above invites a purely geometric bound, which holds regardless of
kinetics.

\begin{theorem}[Footprint bound]
\label{thm:footprint}
Let $\sigma>0$ be the number of precursor molecules deposited per unit area and
let $A>0$ be the area each occupies, with no overlap. Then $\sigma A\le1$, and
if $A_0\le A$ for some ligand footprint $A_0>0$ then
\[
\sigma \le \frac{1}{A_0}.
\]
Equivalently, with growth per cycle $s = p\,\sigma A_{\mathrm{site}}$,
\[
s \le \frac{p\,A_{\mathrm{site}}}{A_0}.
\]
\end{theorem}
\begin{proof}
Non-overlap gives $\sigma A\le1$. Dividing by $\sigma>0$ yields
$A\le1/\sigma$; since $A_0\le A$ and both are positive,
$1/A\le1/A_0$, hence $\sigma\le1/A\le1/A_0$. Multiplying
$s=p\,\sigma A_{\mathrm{site}}$ by this bound gives the second inequality.
\end{proof}

The theorem is a cap, not a prediction: it constrains what a cycle can deposit
and says nothing about whether the cycle reaches the cap. Read backwards,
however, it is a design rule. A ligand of footprint $A_0$ caps the growth per
cycle at $pA_{\mathrm{site}}/A_0$, so a bulky ligand cannot produce a high
growth rate no matter how favourable its surface chemistry.

Read backwards a second time, it becomes testable, because $A_0$ can be
measured independently of any growth experiment. The crystal structure of the
dimer \citep{Kim2016} gives a volume per dimer of
$V=3754/4=938.5$~\AA$^3$, implying a crystal density of
$MZ/(N_A V_{\mathrm{cell}})=1.042$~g\,cm$^{-3}$ --- a value ordinary for an
organometallic, which is a useful check that the cell contents were read
correctly. Treating one dimer as a sphere of that volume gives a cross-section
of
\[
A_0 = \pi\left(\frac{3V}{4\pi}\right)^{2/3} = 115.9\ \text{\AA}^2
  \quad (\text{the cube-root estimate } V^{2/3} \text{ gives } 95.9\ \text{\AA}^2),
\]
against the $A_0\le101.9$--$135.8$~\AA$^2$ that the measured growth per cycle
demands. The two independent published datasets agree, both placing the dimer
footprint at about nine times the bare Ca site area of $12.92$~\AA$^2$.

\begin{remark}
This is consistency, not proof, and the distinction is worth stating precisely.
Agreement of two numbers is evidence that the intact dimer is the adsorbing
species and blocks roughly eight to ten Ca sites per cycle. It does not exclude
a kinetic limitation, which would produce the same measured growth; the crystal
structure is measured at 100~K while deposition runs at 175--250~$^\circ$C, so
the footprint could differ; and the model reaction below is computed for a
truncated ligand. What the comparison does do is convert a vague appeal to
``steric hindrance'' into a falsifiable numerical statement: if a future
measurement placed the adsorbed footprint below about $100$~\AA$^2$, steric
self-limitation by the intact dimer would be excluded.
\end{remark}

\subsection{Why sampling cannot certify a barrier}

It is tempting to sample energies at a set of candidate configurations and read
a barrier off their maximum. That inference is not available, and the failure
is elementary.

\begin{theorem}[No upper bound from sampling]
\label{thm:nosampling}
There is a configuration set $X$ with a subset $S\subseteq X$, nonempty, and an
energy function $E:X\to\R$ such that $E=0$ on all of $S$ while
$\max_X E = 10$. Consequently no function of the sampled values
$\{E(x):x\in S\}$ alone bounds $\max_X E$ from above.
\end{theorem}
\begin{proof}
Take $X=\{0,1,2\}$, $S=\{0,2\}$, and $E(0)=E(2)=0$, $E(1)=10$. Then
$\max_S E = 0$ and $\max_X E = 10$, and the sampled data are identical for this
$E$ and for the function that is identically zero, which has maximum $0$. Any
claimed upper bound derived from the samples would therefore have to be at
least $10$ and at most $0$.
\end{proof}

The positive direction survives: a sampled lower bound is valid at every
configuration where it is evaluated, so its maximum over the samples is a valid
lower bound on any barrier whose path passes through those configurations
(Lemma~\ref{lem:sample-lower}).

\begin{lemma}[Lower bound from samples]
\label{lem:sample-lower}
Let $L(x)\le E(x)$ for $x\in S$ and let $E_r\le U_r$. Then for every $x\in S$,
$L(x)-U_r \le E(x)-E_r$.
\end{lemma}
\begin{proof}
Immediate from $L(x)\le E(x)$ and $-U_r\le-E_r$.
\end{proof}

The asymmetry is the reason a certificate must treat the two directions
separately. When the interval produced for a barrier is one-sided, it must be
reported as one-sided.

\subsection{Formal verification}

Lemmas~\ref{lem:coverage-lean} and \ref{lem:sample-lower} and
Theorems~\ref{thm:footprint} and \ref{thm:nosampling} are formalized in Lean~4
against \texttt{mathlib} as \texttt{gpc\_coverage\_interval},
\texttt{barrier\_lower\_from\_samples},
\texttt{max\_density\_from\_ligand\_size} with
\texttt{gpc\_upper\_from\_ligand\_footprint}, and
\texttt{sampling\_gives\_no\_upper\_bound} respectively, together with the
supporting interval lemmas \texttt{adsorption\_energy\_interval},
\texttt{two\_corrections\_widen}, \texttt{coverage\_gpc\_interval},
\texttt{packing\_bound}, \texttt{barrier\_upper\_at\_max},
\texttt{certified\_barrier\_interval}, \texttt{stays\_positive} and
\texttt{certified\_negative}. Each is checked with \texttt{\#print axioms} and
depends only on \texttt{propext}, \texttt{Classical.choice} and
\texttt{Quot.sound}; no \texttt{sorry}, \texttt{admit} or project-defined axiom
appears in the development. Kernel acceptance certifies the inequalities above
given their hypotheses. It certifies nothing about the hypotheses themselves,
which remain physical or geometric inputs.

\subsection{What remains open for this system}

\begin{enumerate}
\item The elementary-step mechanism and any activation energy for
HF reacting with an adsorbed calcium amidinate on $\mathrm{CaF_2}(111)$. No
value is available in the sources consulted, and Theorem~\ref{thm:nosampling}
shows that a barrier can be certified from below only.
\item The ligand footprint $A_0$ of di-\emph{i}-propylformamidinate, which
Theorem~\ref{thm:footprint} needs in order to become a quantitative cap.
\item Whether the observed $9.51$--$12.68\%$ coverage reflects steric
self-limitation or kinetics. Distinguishing them requires a coverage-dependent
measurement, not a bound.
\item The substrate relationship, which is deliberately absent here. If growth
continues on existing $\mathrm{CaF_2}$ the relevant surface is $\mathrm{CaF_2}$;
the role of an underlying Si(100) is epitaxial orientation, not surface
chemistry.
\end{enumerate}

The honest summary of this section is therefore narrow. Two published numbers
plus a lattice constant fix the fractional coverage of a cycle and show that
site availability is not the limiter. A kernel-checked inequality turns that
observation into a ligand-design cap, and a kernel-checked counterexample shows
that the barrier for the same system cannot be bounded above by any amount of
configuration sampling. No barrier value for this system is claimed.


\section{Evidence and formal verification status}
\label{sec:formalization}
The proofs in this manuscript are mathematical arguments over real and
complex finite-dimensional spaces. They should be distinguished from the
repository's numerical checks and any theorem-prover implementation.
The pairing check builds a fermionic Hamiltonian by sequential application
of creation and annihilation operators. An independent test constructs
Jordan--Wigner tensor-product matrices on the full four-mode Fock space,
checks the canonical anticommutation relations, and projects to the explicit
two-electron sector. Comparing these constructions can detect occupation-order
signs, omitted self-pairing terms, and an incorrect energy origin. Tests of
normalized vectors, both paired eigenvalues, the complete sector spectrum, and
degenerate cases are stronger than comparison of one energy alone. They remain
tests of a toy model in floating-point arithmetic.

The dense block experiment evaluates small finite matrices and checks their
enclosures against direct diagonalization. It does not demonstrate efficient
complement certification or validated numerical rounding. The regression
checks in \path{tests/test_audit_counterexamples.py} have also passed: they
reproduce the distinction between component count and independence number,
the identical 2-RDMs but different 3-RDMs in Example~\ref{ex:three-rdm},
and the exclusion of a physical state by the invalid proposed SOS constraint.
These are negative validation tests of mathematical claims, not chemical data
or tests of a surface-reaction model.

\begin{table}[ht]
\centering
\caption{Reproduced finite-model checks. All numerical references are ordinary
floating-point computations; none is a molecular or experimental benchmark.}
\begin{tabular}{lll}
\hline
Check & Cases & Observed result \\
\hline
Python regression suite & 30 tests & All passed \\
Random Hermitian block splits & 360 & Max. interval violation $4.44\times10^{-16}$ \\
Sharp scalar-block cases & 9 & Max. endpoint error $8.88\times10^{-16}$ \\
Finite CTMC comparisons & 20 & No violation above tested tolerance \\
Pairing CLI parameter cases & 10 & All passed \\
\hline
\end{tabular}
\label{tab:reproduced}
\end{table}
The block experiment uses dimensions two through ten, real and complex
entries, four samples per dimension and kind, every nontrivial leading
split, and seed 20261112. Acceptance tolerances, library versions, every
case, and the synthetic kinetic parameters are stored in
\path{experiments/certificate_results.json}. A nonzero rounding violation is
reported rather than silently replaced by an exact enclosure.

Historical hard-coded SOS energy tables and hand-assigned surface reference
counts are excluded from the evidence base. No physical accuracy, speedup, or
independent DFT validation is inferred from them.

\paragraph{Lean source mapping and build boundary.}
The current development uses real numbers and mathlib v4.19.0. It has been
compiled with Lean 4.19.0 and every declaration listed below was checked with
\texttt{\#print axioms}: each depends only on the three standard Lean axioms
\texttt{propext}, \texttt{Classical.choice} and \texttt{Quot.sound}. No
\texttt{sorry}, no \texttt{admit} and no project-specific axiom occurs in the
verified closure. Reproduce with \texttt{lake build} followed by
\texttt{lake env lean scripts/CheckAxioms.lean}. Kernel acceptance certifies the
scalar lemmas that are listed, not the physical modelling steps discussed
below. Unless qualified otherwise, the names below are in the
\texttt{CalcChemistry} namespace in
\path{CalcChemistry/CertifiedTheory.lean}.
\begin{sloppypar}
\begin{itemize}
\item \path{reaction_interval} proves real interval subtraction as in
Proposition~\ref{prop:barrier}. \path{add_error_interval} propagates an
additive correction whose absolute value is bounded by a supplied radius.
It assumes the correction bound; it does not derive a physical model error.
\item \path{psd_quadratic_nonneg} is a real two-variable quadratic-form
inequality under nonnegative diagonal and determinant conditions.
\path{sharp_block_scalar_lower} proves the sharp square-root bound for the
scalar quadratic form used in Theorem~\ref{thm:block}. It does not construct
matrix blocks or prove their spectral hypotheses.
\item \path{quadratic_gap_lower} proves the scalar inequality with
$a-b^2/(c-a)$ under $c-a>0$. \path{normalized_block_lower} transfers it
to a supplied energy using a normalization identity and a supplied
quadratic-form lower bound, as in Corollary~\ref{cor:gap}.
\item \path{ground_energy_interval} transfers a uniform lower bound and
one trial upper bound to an abstract attained Rayleigh minimum. The existence
and attainment of that minimum are explicit assumptions. The identification
with a Hermitian matrix's least eigenvalue is not formalized by this theorem.
\item \path{exponential_barrier_interval} proves the real exponential
monotonicity used in Theorem~\ref{thm:rate}; \path{rate_interval} also
allows an uncertain nonnegative prefactor. The strictly positive prefactor
hypotheses in this manuscript support the further logarithmic formulas.
Neither declaration proves the transition-state rate law itself.
\item \path{relaxation_lower_bound} proves the order implication from an
outer inclusion, a supplied lower bound throughout the relaxed set, and an
attained physical energy. It does not implement an RDM SDP or verify a dual
matrix slack. \path{certified_ranking} proves strict ordering from disjoint
sound real intervals; the kinetic interpretation requires the rate intervals
proved separately.
\item \path{pairing_characteristic} proves the paired-block polynomial
identity and \path{pairing_lower_root} proves that the lower square-root
expression is a root. These declarations do not construct the CAR operators,
identify the full six-state spectrum, or establish a molecular model.
\item \path{BornOppenheimer.parameter_nonneg}, in
\path{CalcChemistry/BornOppenheimer.lean}, proves nonnegativity of the mass
parameter under the stated mass hypotheses. It provides no BO error estimate.
\item \path{block_energy_certificate}, in
\path{CalcChemistry/BarrierBounds.lean}, lifts the scalar inequality above to
the block operator in arbitrary finite dimensions $n$ and $m$: given
Rayleigh-quotient lower bounds $a$ and $c$ on the two diagonal blocks and a
bound $b$ on the coupling form, every state obeys the sharp two-block inequality
with constant $(a+c-\sqrt{(a-c)^2+4b^2})/2$. \path{block_energy_certificate_gap}
gives the same bound in the square-root-free form $a-b^2/(c-a)$ for $c>a$, which
is the downfolding correction quoted in practice, and
\path{normalized_energy_lower} states it for a unit vector. The block
hypotheses remain explicit assumptions: nothing here constructs the blocks from
molecular integrals or certifies the omitted space.
\item The surface-reaction certificate of Section~\ref{sec:worked} is
formalized in \path{CalcChemistry/SurfaceReaction.lean}:
\path{gpc_coverage_interval} and \path{coverage_gpc_interval} prove the interval
relation between growth per cycle, layer spacing and fractional coverage;
\path{adsorption_energy_interval} and \path{two_corrections_widen} propagate
intervals through an energy difference and through additive corrections;
\path{packing_bound}, \path{max_density_from_ligand_size} and
\path{gpc_upper_from_ligand_footprint} prove the ligand-footprint design rule of
Theorem~\ref{thm:footprint}, and \path{steric_admissible} and
\path{steric_falsification} state it as an admissibility test together with the
condition that would refute steric control.
\path{barrier_lower_from_samples}, \path{barrier_upper_at_max} and
\path{certified_barrier_interval} give the two-sided barrier interval, and
\path{sampling_gives_no_upper_bound} is the counterexample of
Theorem~\ref{thm:nosampling} proving that a finite sample cannot bound a barrier
from above. \path{stays_positive} and \path{certified_negative} convert
intervals into sign statements.
\end{itemize}
\end{sloppypar}

This gives $33$ kernel-checked declarations in total, every one of them
audited with \texttt{\#print axioms}. The count is of declarations, not of
scientific content: all but the block certificate and the ligand inequalities
are real-arithmetic plumbing, and none of them certifies a molecular energy.

The energy and fifth-order certificate structures store soundness inequalities
as fields supplied by the caller. Storing such an inequality does not prove
it for a physical system. The finite-dimensional matrix-to-Rayleigh argument,
the CTMC semigroup theorem and Duhamel proof, and a physical BO approximation
theorem are not claimed as formalized here; they are proved on paper in the
preceding sections and remain outside the kernel-checked closure. The exact
declarations, the pinned Lean and mathlib versions, and the reproducible build
that produced the axiom report are recorded in the repository status document.

\section{Conclusion}
For a specified finite Hamiltonian, a justified complement bound and coupling
bound turn an active calculation into a two-sided energy certificate.
Variational trial states and dual-feasible relaxations supply complementary
information. Barrier subtraction, monotone rate propagation, and finite-state
Markov stability then give conditional conclusions with explicit dependence on
energy precision, prefactor uncertainty, elapsed time, and the event catalog.
The unresolved scientific task is to produce these inputs for a defined
surface-reaction model and to bound the changes between that model and the
physical target. A wide interval or an unverified hypothesis is an informative
outcome: it identifies the missing evidence without claiming a kinetic
prediction that the calculation cannot support.

\begin{thebibliography}{22}
\bibitem[O'Gorman et~al.(2022)O'Gorman, Irani, Whitfield, and
  Fefferman]{OGorman2022}
Bryan O'Gorman, Sandy Irani, James Whitfield, and Bill Fefferman.
\newblock Intractability of electronic structure in a fixed basis.
\newblock \emph{PRX Quantum}, 3\penalty0 (2):\penalty0 020322, 2022.
\newblock \doi{10.1103/PRXQuantum.3.020322}.
\newblock URL \url{https://doi.org/10.1103/PRXQuantum.3.020322}.
\bibitem[Liu et~al.(2007)Liu, Christandl, and Verstraete]{Liu2007}
Yi-Kai Liu, Matthias Christandl, and F.~Verstraete.
\newblock Quantum computational complexity of the {N}-representability problem:
  {QMA} complete.
\newblock \emph{Physical Review Letters}, 98\penalty0 (11):\penalty0 110503,
  2007.
\newblock \doi{10.1103/PhysRevLett.98.110503}.
\newblock URL \url{https://doi.org/10.1103/PhysRevLett.98.110503}.
\bibitem[Pieck and Tonner-Zech(2025)]{Pieck2025}
Fabian Pieck and Ralf Tonner-Zech.
\newblock Computational ab initio approaches for area-selective atomic layer
  deposition: Methods, status, and perspectives.
\newblock \emph{Chemistry of Materials}, 37\penalty0 (9):\penalty0 2979--3021,
  2025.
\newblock \doi{10.1021/acs.chemmater.4c03477}.
\newblock URL \url{https://doi.org/10.1021/acs.chemmater.4c03477}.
\bibitem[Mazziotti(2012)]{Mazziotti2012}
David~A. Mazziotti.
\newblock Two-electron reduced density matrix as the basic variable in
  many-electron quantum chemistry and physics.
\newblock \emph{Chemical Reviews}, 112\penalty0 (1):\penalty0 244--262, 2012.
\newblock \doi{10.1021/cr2000493}.
\newblock URL \url{https://doi.org/10.1021/cr2000493}.
\bibitem[Richardson(1963)]{Richardson1963}
R.~W. Richardson.
\newblock A restricted class of exact eigenstates of the pairing-force
  {Hamiltonian}.
\newblock \emph{Physics Letters}, 3\penalty0 (6):\penalty0 277--279, 1963.
\newblock \doi{10.1016/0031-9163(63)90259-2}.
\newblock URL \url{https://doi.org/10.1016/0031-9163(63)90259-2}.
\bibitem[Dukelsky et~al.(2004)Dukelsky, Pittel, and Sierra]{Dukelsky2004}
J.~Dukelsky, S.~Pittel, and G.~Sierra.
\newblock Colloquium: Exactly solvable {Richardson-Gaudin} models for many-body
  quantum systems.
\newblock \emph{Reviews of Modern Physics}, 76\penalty0 (3):\penalty0 643--662,
  2004.
\newblock \doi{10.1103/RevModPhys.76.643}.
\newblock URL \url{https://doi.org/10.1103/RevModPhys.76.643}.
\bibitem[Hagedorn(1988)]{Hagedorn1988}
George~A. Hagedorn.
\newblock High order corrections to the time-independent {Born-Oppenheimer}
  approximation {II}: Diatomic {Coulomb} systems.
\newblock \emph{Communications in Mathematical Physics}, 116\penalty0
  (1):\penalty0 23--44, 1988.
\newblock \doi{10.1007/BF01239023}.
\newblock URL \url{https://doi.org/10.1007/BF01239023}.
\bibitem[Eyring(1935)]{Eyring1935}
Henry Eyring.
\newblock The activated complex in chemical reactions.
\newblock \emph{The Journal of Chemical Physics}, 3\penalty0 (2):\penalty0
  107--115, 1935.
\newblock \doi{10.1063/1.1749604}.
\newblock URL \url{https://doi.org/10.1063/1.1749604}.
\bibitem[Vogt et~al.(2005)Vogt, Henning, and Weiss]{Vogt2005}
J.~Vogt, J.~Henning, and H.~Weiss.
\newblock The structure of caf2(111) and baf2(111) single crystal surfaces: A
  tensor low energy electron diffraction study.
\newblock \emph{Surface Science}, 578\penalty0 (1-3):\penalty0 57--70, 2005.
\newblock \doi{10.1016/j.susc.2005.01.011}.
\newblock URL \url{https://doi.org/10.1016/j.susc.2005.01.011}.
\bibitem[Puchin et~al.(2001)Puchin, Puchina, Huisinga, and
  Reichling]{Puchin2001}
V~E Puchin, A~V Puchina, M~Huisinga, and M~Reichling.
\newblock Theoretical modelling of steps on the caf2(111) surface.
\newblock \emph{Journal of Physics: Condensed Matter}, 13\penalty0
  (10):\penalty0 2081--2094, 2001.
\newblock \doi{10.1088/0953-8984/13/10/304}.
\newblock URL \url{https://doi.org/10.1088/0953-8984/13/10/304}.
\bibitem[Calleja et~al.(2005)Calleja, Hinarejos, Vázquez~de Parga, Suturin,
  Sokolov, and Miranda]{Calleja2005}
F.~Calleja, J.J. Hinarejos, A.L. Vázquez~de Parga, S.M. Suturin, N.S. Sokolov,
  and R.~Miranda.
\newblock Epitaxial growth of caf2(111) on cu(111) visualized by {STM}.
\newblock \emph{Surface Science}, 582\penalty0 (1-3):\penalty0 14--20, 2005.
\newblock \doi{10.1016/j.susc.2005.03.003}.
\newblock URL \url{https://doi.org/10.1016/j.susc.2005.03.003}.
\bibitem[Hennessy et~al.(2023)Hennessy, Rodríguez, West, and
  Brandon]{Hennessy2023}
John Hennessy, Robin~E. Rodríguez, William~C. West, and Erik~J. Brandon.
\newblock Atomic layer deposition of calcium fluoride for barrier coating
  applications.
\newblock \emph{Journal of Vacuum Science \& Technology A}, 41\penalty0
  (4):\penalty0 042405, 2023.
\newblock \doi{10.1116/6.0002563}.
\newblock URL \url{https://doi.org/10.1116/6.0002563}.
\bibitem[Kim et~al.(2016)Kim, Yang, Powers, Davis, Lou, and Gordon]{Kim2016}
Sang~Bok Kim, Chuanxi Yang, Tamara Powers, Luke~M. Davis, Xiabing Lou, and
  Roy~G. Gordon.
\newblock Synthesis of calcium({II}) amidinate precursors for atomic layer
  deposition through a redox reaction between calcium and amidines.
\newblock \emph{Angewandte Chemie International Edition}, 55\penalty0
  (35):\penalty0 10228--10233, 2016.
\newblock \doi{10.1002/anie.201602406}.
\newblock URL \url{https://doi.org/10.1002/anie.201602406}.
\bibitem[Kaawar and Paulus(2019)]{Kaawar2019}
Zeinab Kaawar and Beate Paulus.
\newblock Adsorption of hydrogen fluoride on alkaline earth fluoride surfaces:
  A first-principles study.
\newblock \emph{Journal of Fluorine Chemistry}, 224:\penalty0 67--72, 2019.
\newblock \doi{10.1016/j.jfluchem.2018.09.006}.
\newblock URL \url{https://doi.org/10.1016/j.jfluchem.2018.09.006}.
\bibitem[Pilvi et~al.(2007)Pilvi, Arstila, Leskelä, and Ritala]{Pilvi2007}
Tero Pilvi, Kai Arstila, Markku Leskelä, and Mikko Ritala.
\newblock Novel {ALD} process for depositing caf 2 thin films.
\newblock \emph{Chemistry of Materials}, 19\penalty0 (14):\penalty0 3387--3392,
  2007.
\newblock \doi{10.1021/cm0629412}.
\newblock URL \url{https://doi.org/10.1021/cm0629412}.
\bibitem[Pilvi et~al.(2008)Pilvi, Ritala, Leskelä, Bischoff, Kaiser, and
  Kaiser]{Pilvi2008}
Tero Pilvi, Mikko Ritala, Markku Leskelä, Martin Bischoff, Ute Kaiser, and
  Norbert Kaiser.
\newblock Atomic layer deposition process with tif\_4 as a precursor for
  depositing metal fluoride thin films.
\newblock \emph{Applied Optics}, 47\penalty0 (13):\penalty0 C271, 2008.
\newblock \doi{10.1364/AO.47.00C271}.
\newblock URL \url{https://doi.org/10.1364/AO.47.00C271}.
\bibitem[Lee et~al.(2016)Lee, Sun, Young, and George]{Lee2016}
Younghee Lee, Huaxing Sun, Matthias~J. Young, and Steven~M. George.
\newblock Atomic layer deposition of metal fluorides using {HF}–pyridine as
  the fluorine precursor.
\newblock \emph{Chemistry of Materials}, 28\penalty0 (7):\penalty0 2022--2032,
  2016.
\newblock \doi{10.1021/acs.chemmater.5b04360}.
\newblock URL \url{https://doi.org/10.1021/acs.chemmater.5b04360}.
\bibitem[Lo~Presti et~al.(2022)Lo~Presti, Pellegrino, and
  Malandrino]{LoPresti2022}
Francesca Lo~Presti, Anna~L. Pellegrino, and Graziella Malandrino.
\newblock Journey of a molecule from the solid to the gas phase and vice versa
  : direct estimation of vapor pressure of alkaline-earth metalorganic
  precursors for atmospheric pressure vapor phase deposition of fluoride films.
\newblock \emph{Dalton Transactions}, 51\penalty0 (18):\penalty0 7352--7362,
  2022.
\newblock \doi{10.1039/d2dt00479h}.
\newblock URL \url{https://doi.org/10.1039/d2dt00479h}.
\bibitem[Ylilammi and Ranta‐aho(1994)]{Ylilammi1994}
M.~Ylilammi and T.~Ranta‐aho.
\newblock Metal fluoride thin films prepared by atomic layer deposition.
\newblock \emph{Journal of The Electrochemical Society}, 141\penalty0
  (5):\penalty0 1278--1284, 1994.
\newblock \doi{10.1149/1.2054910}.
\newblock URL \url{https://doi.org/10.1149/1.2054910}.
\bibitem[Lee et~al.(2015)Lee, DuMont, Cavanagh, and George]{Lee2015}
Younghee Lee, Jaime~W. DuMont, Andrew~S. Cavanagh, and Steven~M. George.
\newblock Atomic layer deposition of alf3 using trimethylaluminum and hydrogen
  fluoride.
\newblock \emph{The Journal of Physical Chemistry C}, 119\penalty0
  (25):\penalty0 14185--14194, 2015.
\newblock \doi{10.1021/acs.jpcc.5b02625}.
\newblock URL \url{https://doi.org/10.1021/acs.jpcc.5b02625}.
\bibitem[Ma et~al.(2011)Ma, Guo, Gordon, and Zaera]{Ma2011}
Qiang Ma, Hansheng Guo, Roy~G. Gordon, and Francisco Zaera.
\newblock Surface chemistry of copper(i) acetamidinates in connection with
  atomic layer deposition ({ALD}) processes.
\newblock \emph{Chemistry of Materials}, 23\penalty0 (14):\penalty0 3325--3334,
  2011.
\newblock \doi{10.1021/cm200432t}.
\newblock URL \url{https://doi.org/10.1021/cm200432t}.
\bibitem[Khumaini et~al.(2025)Khumaini, Cho, Kim, and Lee]{Khumaini2025}
Khabib Khumaini, Gyejun Cho, Hye-Lee Kim, and Won-Jun Lee.
\newblock Thermal atomic layer etching mechanism of amorphous aluminum oxide by
  hydrogen fluoride and trimethylaluminum: First-principles study.
\newblock \emph{Surfaces and Interfaces}, 72:\penalty0 107114, 2025.
\newblock \doi{10.1016/j.surfin.2025.107114}.
\newblock URL \url{https://doi.org/10.1016/j.surfin.2025.107114}.
\end{thebibliography}
\end{document}
